\section{Referencial Teórico}

\subsection{Teoria do Fluxo}
A Teoria do Fluxo, proposta por Mihaly Csikszentmihalyi, descreve um estado psicológico de imersão total e envolvimento intenso durante a realização de uma atividade. Seu princípio central é o equilíbrio entre o nível de habilidade do indivíduo e o desafio imposto pela tarefa: quando esse equilíbrio é atingido, o indivíduo entra no estado de fluxo, caracterizado por foco intenso, motivação intrínseca, sensação de controle e uma experiência autotélica, sendo assim a realização da atividade pelo prazer em si. Em contrapartida, tarefas muito fáceis tendem a gerar tédio, enquanto tarefas excessivamente difíceis provocam ansiedade. Do ponto de vista neurológico, o estado de fluxo eleva a cognição, otimizando a concentração, a velocidade de processamento e a tomada de decisões \cite{csikszentmihalyi1989optimal}.

\subsection{Teoria da Autodeterminação (SDT)}
A Teoria da Autodeterminação, desenvolvida por Deci e Ryan, é uma teoria da motivação humana que se distingue por analisar não apenas a quantidade de motivação de um indivíduo, mas sobretudo, a sua qualidade. Apresentada como uma abordagem dialética, a teoria investiga as necessidades psicológicas básicas de autonomia, competência e pertencimento, argumentando que essas necessidades são inatas a todos os indivíduos \cite{cernev2012autodeterminacao}. A autonomia refere-se ao desejo de agir por vontade própria, orientado por interesses e valores internos, e não por pressão externa. A competência diz respeito à necessidade de dominar desafios em um nível adequado às próprias habilidades e de receber \textit{feedback} positivo sobre o desempenho. Já o pertencimento expressa o desejo de se sentir valorizado e conectado a um grupo, por meio de vínculos sociais genuínos.\\
\indent A satisfação dessas necessidades, contudo, não depende exclusivamente do indivíduo, ela está diretamente condicionada ao ambiente em que ele está inserido. Os fatores ambientais podem facilitar ou prejudicar o sentimento de autodeterminação, contribuindo para a qualidade do desempenho e a sensação de bem-estar pessoal.

\subsection{Ajuste Dinâmico de Dificuldade (DDA)}
O Ajuste Dinâmico de Dificuldade é uma técnica que modifica automaticamente parâmetros de um sistema em tempo real com o objetivo de manter o usuário em um nível ótimo de desafio. Para operar, um sistema DDA é composto por dois elementos centrais: uma métrica de dificuldade, responsável por quantificar numericamente o desafio percebido pelo usuário e um mecanismo de ajuste, que define como os parâmetros do sistema são alterados para manter o equilíbrio entre desafio e proficiência \cite{oliveira2025dda}. 

\subsection{Trabalhos relacionados}
A literatura de gamificação na saúde digital mostra que abordagens uniformes de engajamento tendem a perder eficácia quando aplicadas a perfis diferentes de usuários. Por exemplo, Plangger et al. (2022) apontam que recompensas simples e métricas padronizadas apresentam impacto limitado no comportamento de longo prazo, porque não avaliam adequadamente o nível de condicionamento e a experiência subjetiva do usuário \cite{plangger2022little}. Estas evidências reforçam que a motivação sustentável depende de adaptação contínua ao estado individual de quem treina.

Este trabalho foi inspirado no estudo de Bastos et al. (2023), que investigou um serious game baseado em gamificação e biofeedback para aumentar a adesão à atividade física \cite{bastos2023gamification}. A pesquisa de Bastos et al. serviu de referência ao demonstrar a viabilidade de integrar sinais biométricos em uma experiência lúdica, mas também revelou que a maioria das soluções ainda não explora de forma robusta a adaptação do desafio em tempo real.

Na prática de mercado, plataformas como \textit{Fitbit}, \textit{Apple Fitness+} e \textit{Garmin} ainda oferecem principalmente feedback reativo e genérico, focado em metas batidas, contagem de passos e calorias acumuladas. Há uma lacuna clara para abordagens capazes de fornecer feedback preditivo e personalizado, que indiquem se o usuário está abaixo, dentro ou acima do canal de fluxo durante a sessão, e não apenas ao término dela. A identificação desta lacuna motivou nossa proposta e reforçou a importância de buscar uma solução que combine biofeedback, gamificação e ajuste dinâmico de dificuldade para manter a motivação e a retenção do usuário.